\chapter{The Statement of Work and Project Plan}

\section{The Statement of Work}

\begin{tipbox}
For the sprint-level checklist of SOW deliverables and deadlines, see the \textbf{Student Guide, Sprint~1 section}. This chapter provides the methodology; the Student Guide provides the week-by-week action items.
\end{tipbox}

The Statement of Work (SOW) is the foundational agreement between your team, the PA, and the client. It defines what will be done, by when, within what constraints, and by whom. The SOW is developed during SDI Sprint~1 (Weeks~3--4) and submitted as a formal deliverable.

The SOW is not a formality. It is the document that prevents scope creep, resolves misunderstandings about deliverables, and protects the team when the client asks for something that was never agreed upon. In professional practice, the SOW is legally binding. In capstone, it is academically binding: your PA evaluates your project against what you committed to deliver.

\subsection{SOW Structure}
A complete SOW contains the following sections:

\begin{enumerate}[nosep,leftmargin=1.5em]
\item \textbf{Project Title and Team}: project name, team number, all members with assigned roles (Scrum Master, CFO, Technical Lead, Communication Lead, Fabrication Lead), PA name and contact, client name and contact, Technical Advisor if identified.
\item \textbf{Track Classification}: Build, Paper, or Competition with justification. Determines PoC expectations and final deliverable type.
\item \textbf{Problem Statement}: Five-Element Problem Statement from Chapter~2.
\item \textbf{Scope}: bounded description of what is included and what is \textbf{explicitly excluded}. The ``Out of Scope'' section is as important as the ``In Scope'' section. List features the client mentioned that will not be addressed with a rationale for each exclusion.
\item \textbf{Deliverables}: tangible outputs per semester. SDI: SOW, SDRD, PDR report and presentation, PoC. SDII: CDR report and presentation, FDR report and presentation, tested prototype or technical package, O\&M documentation, Showcase poster.
\item \textbf{Schedule}: milestones mapped to sprints with dates. Include client meetings, design reviews, procurement deadlines, fabrication milestones, test milestones, and Showcase.
\item \textbf{Budget Estimate}: major categories (materials, components, services, travel, contingency) with funding source.
\item \textbf{Team Roles and Responsibilities}: who is responsible for each major activity with specific duties.
\item \textbf{Assumptions and Constraints}: what you assume to be true (client resources, facility access, technology availability) and what bounds the design (budget, regulations, physical space).
\item \textbf{Risk Summary}: top 3--5 risks with likelihood, impact, and preliminary mitigation.
\end{enumerate}

\subsection{Track Classification}
Track classification is documented in the SOW and approved by the PA, client, and CF. It \textbf{cannot be changed after PDR} without written approval from the Capstone Program Director.

\subsection{Scope Change Management}
When the client requests a change to the agreed scope (and they will), follow this process:
\begin{enumerate}[nosep,leftmargin=1.5em]
\item \textbf{Document the request}: write down exactly what is being asked. Confirm understanding with the client.
\item \textbf{Assess impact}: how does this affect schedule, budget, requirements, and risk?
\item \textbf{Present the trade-off}: show the client what must be deferred or dropped to accommodate the change.
\item \textbf{Get explicit approval}: client and PA must both agree. Document the decision in the SOW.
\end{enumerate}

Without this discipline, small changes accumulate until the project bears no resemblance to the original agreement.

\begin{warningbox}[The SOW Is a Scope Control Document]
When a client asks for additional features, the SOW is your reference: ``That feature is not in the current scope. We can add it, but here is the impact on schedule and budget. Do you want to proceed?'' Without a baselined SOW, every client request becomes an obligation and scope creep consumes the semester.
\end{warningbox}

\section{The Initial Project Plan}
Developed during Sprint~1 (Weeks~3--4), the Project Plan expands the SOW into an operational document:
\begin{itemize}[nosep,leftmargin=1.5em]
\item \textbf{Product Backlog}: all known work items, prioritized by client value (MoSCoW).
\item \textbf{First Sprint Backlog}: items selected for Sprint~1 with owners and effort estimates.
\item \textbf{Team norms}: meeting schedule, communication tools, response time expectations, file conventions, Team Contract.
\item \textbf{Risk register}: expanded from SOW with mitigations assigned to specific owners.
\end{itemize}

The Plan is emailed to the client and uploaded to Canvas. Reviewed and updated at every retrospective.

\section{Project Goals Setting}
At each semester's start, set SMART goals with PA and client agreement: Specific, Measurable, Achievable, Relevant, Time-bound. Examples:
\begin{itemize}[nosep,leftmargin=1.5em]
\item ``Demonstrate a working temperature control loop with $\pm$2\,\degC{} accuracy by PDR'' (SDI)
\item ``Complete a validated FEA model under ASCE~7 load cases by CDR'' (SDII)
\item ``Achieve $\geq$85\% of VCRM requirements verified as PASS by FDR'' (SDII)
\end{itemize}

At semester's end, the \textbf{Project Goal Attainment} document evaluates performance. Teams not accomplishing goals will not receive an A.

\section{Scope Management Throughout the Project}
Scope management is the discipline of controlling what is and is not included in the project. It is one of the most important project management skills and one of the hardest to maintain under the pressure of client enthusiasm and team ambition.

\subsection{The Scope Triangle}
Every project is constrained by three interdependent variables: \textbf{scope} (what the system does), \textbf{schedule} (when it is delivered), and \textbf{budget} (what it costs). Changing one variable affects the others. If the client adds a feature (increasing scope), either the schedule extends, the budget increases, or both. If the schedule compresses (e.g., a competition date moves earlier), either scope decreases or budget increases (e.g., expedited shipping).

Your SOW fixes the scope-schedule-budget triangle at a specific equilibrium. Every scope change request must be evaluated against this triangle: what is the impact on schedule and budget? If the impact is acceptable to all parties (team, PA, client), proceed with the change under documented change control. If not, the change is deferred to future work.

\subsection{Writing Effective Deliverable Descriptions}
Each deliverable in the SOW must be described with enough specificity that all parties agree on what ``done'' looks like. Bad: ``Provide a final report.'' Good: ``Deliver a Final Design Report (FDRt) containing: executive summary, complete VCRM with all requirements dispositioned, as-designed vs. as-built comparison, test data and analysis, O\&M documentation, lessons learned, and future work recommendations. Format per the Capstone Design Program Format Guide. Submitted to PA and client no later than [date].''

The deliverable description is the acceptance criterion. At FDR, if the client says ``I expected more detail in the O\&M section,'' you can point to the SOW: ``The O\&M documentation was specified as [description]. If you need additional detail, we can discuss a scope extension.''

\subsection{Common SOW Mistakes}
\begin{enumerate}[nosep,leftmargin=1.5em]
\item \textbf{Scope too broad}: ``Design and build a complete autonomous vehicle.'' No capstone team can build a complete autonomous vehicle in two semesters. Narrow the scope to a specific subsystem or capability.
\item \textbf{Scope too vague}: ``Improve the client's process.'' Improve how? By what metric? To what target? Vague scope leads to vague deliverables and ambiguous evaluation.
\item \textbf{Missing ``Out of Scope''}: not explicitly stating what is excluded. The client assumes features are included unless they are explicitly excluded.
\item \textbf{No budget contingency}: budgeting exactly to the penny with no reserve for unexpected costs.
\item \textbf{Unrealistic schedule}: promising deliverables faster than the sprint structure allows.
\item \textbf{Missing risk acknowledgment}: not identifying that the project has known technical risks that could affect scope.
\end{enumerate}

\subsection{SOW as a Professional Development Tool}
Writing a SOW is a career skill. In consulting, the SOW (or proposal) is the document that wins (or loses) the contract. In industry, the SOW defines the work breakdown for your team. In government contracting, the SOW is a legally binding document with financial consequences for non-compliance. The SOW you write in capstone is practice for every project you will manage in your career.


\section{The SOW Review Process}
The SOW should be reviewed by all stakeholders before baselining:

\textbf{Team review}: every member reads the complete SOW and verifies that their role, responsibilities, and the scope description are accurate. Each member signs or acknowledges the SOW.

\textbf{PA review}: the PA checks for: realistic scope (is this achievable in two semesters?), adequate budget (does the estimate cover all foreseeable costs?), appropriate track classification (does the track match the project deliverables?), identified risks (are the top risks realistic and mitigable?), and completeness (are all required sections present?).

\textbf{Client review}: the client verifies that the scope, deliverables, and schedule match their expectations. The client's approval is documented (email confirmation at minimum; signature on the SOW is preferred). After client approval, the SOW is baselined.

\section{SOW Templates and Examples}
The Capstone Design Program provides a SOW template on Canvas. Use it as a starting point but do not treat it as a fill-in-the-blank form. Your SOW should reflect the unique characteristics of your project, not generic placeholder text. Review SOW examples from previous capstone teams (available from your PA or on Canvas) to calibrate the expected level of detail.

\begin{tipbox}
Write the SOW as if you are handing it to a new team that knows nothing about your project. Can they understand the scope, deliverables, and constraints from the SOW alone? If not, add detail. The SOW should be self-contained.
\end{tipbox}



\section{Negotiating Scope with Clients}
Scope negotiation is one of the most important professional skills you practice in capstone. Clients often want more than is achievable in two semesters. Your job is not to say ``no'' but to present options with trade-offs.

\textbf{The ``Yes, and'' technique}: instead of rejecting a request, acknowledge it and present the cost. ``Yes, we can add wireless connectivity, and the impact would be 3 additional weeks of development and approximately \$150 in components. To stay on schedule, we would need to defer the color display to a future version. Which is more important to you?''

\textbf{Minimum Viable Product (MVP)}: identify the smallest set of features that delivers value to the client. This becomes your ``Must Have'' scope. Everything else is ``Should'' or ``Could.'' If the project goes perfectly, you deliver it all. If challenges arise (and they will), the MVP is protected.

\textbf{Version planning}: frame the two-semester capstone as ``Version 1.0.'' Features that cannot fit in V1.0 are documented as ``Future Work'' in the SOW and the FDR report. This respects the client's vision while being honest about what is achievable. Clients who understand version planning are far more satisfied than clients who expect everything and receive a rushed partial implementation.

\textbf{The 80/20 rule}: in most projects, 80\% of the value comes from 20\% of the features. Identify the high-value features early and protect them. The remaining features; the ones that are nice-to-have but do not fundamentally change the system's value proposition; are the first candidates for deferral when schedule pressure mounts.



\section{Building the Risk Register}
The SOW includes a preliminary risk summary (top 3--5 risks). The full \textbf{risk register}, developed during Sprint~1--2 and maintained throughout the project, expands this into a living management tool.

Each risk entry contains: (1)~\textbf{Risk ID} (R-001, R-002, etc.), (2)~\textbf{Description} (specific, not vague: ``Thermocouple accuracy degrades above 800\,\degC{} due to oxidation'' not ``temperature measurement may be inaccurate''), (3)~\textbf{Likelihood} (1--5 scale), (4)~\textbf{Impact} (1--5 scale), (5)~\textbf{Risk Score} (Likelihood $\times$ Impact), (6)~\textbf{Mitigation strategy} (what action reduces the likelihood or impact), (7)~\textbf{Owner} (which team member is responsible for monitoring and executing the mitigation), and (8)~\textbf{Status} (Open, Mitigated, Retired, Occurred).

\textbf{Risk categories to consider}:
\begin{itemize}[nosep,leftmargin=1.5em]
\item \textbf{Technical}: sensor accuracy, actuator performance, software bugs, EMI, thermal limits, structural failure.
\item \textbf{Schedule}: procurement delays, fabrication overruns, testing reveals redesign need, team member availability (illness, other course demands, job interviews).
\item \textbf{Budget}: component costs higher than estimated, shipping surcharges, PCB revision costs, travel costs exceed estimate.
\item \textbf{External}: client changes priorities, regulatory requirement discovered late, facility access restricted (InnoHub closure, equipment breakdown).
\item \textbf{Team}: unequal contribution, key member drops the course, personality conflicts escalate, knowledge concentrated in one person (bus factor = 1).
\end{itemize}

Review the risk register at every sprint retrospective. Update likelihood and status based on new information. A risk register that has not been updated since Sprint~1 is not a management tool; it is a decoration.

\section{Documenting Assumptions and Constraints}
The SOW should explicitly document what you are assuming to be true and what constraints bound the design:

\textbf{Assumptions} are things you believe to be true but have not verified. Examples: ``The client will provide access to the test site by Sprint~4,'' ``The operating temperature will not exceed 50\,\degC{},'' ``Shop air supply provides 90\,psi minimum.'' When an assumption proves false, it becomes a risk or a scope change.

\textbf{Constraints} are non-negotiable boundaries. Examples: ``Total budget not to exceed \$2,500,'' ``System must fit within a 12''$\times$12''$\times$8'' enclosure,'' ``Must comply with IEC~61010 safety standard,'' ``Must use ESP32 microcontroller per client specification.''

Documenting assumptions protects the team: when the client says ``I assumed you would also build a mobile app,'' you can point to the SOW assumptions section and initiate a scope change process rather than absorbing unplanned work.



\section{The Stakeholder Communication Matrix}
Different stakeholders need different information at different frequencies. Document this in a communication matrix:

\begin{center}\small
\begin{tabularx}{\textwidth}{l>{\raggedright\arraybackslash}X>{\raggedright\arraybackslash}X>{\raggedright\arraybackslash}X}
\toprule
\textbf{Stakeholder} & \textbf{Information Need} & \textbf{Frequency} & \textbf{Method} \\
\midrule
Client & Progress, decisions, blockers & Biweekly & Status email + meetings \\
PA & Sprint progress, risks, grades & Weekly & Sprint review + email \\
Technical Advisor & Technical approach, analysis & As needed & Meeting + report sections \\
Capstone Purchasing & Orders, budget status & Per purchase & Forms + email \\
InnoHub Staff & Fabrication plans, training & Per session & In-person + email \\
Team Members & Tasks, blockers, decisions & Daily & Stand-up + task board \\
\bottomrule
\end{tabularx}
\end{center}

The Communication Lead owns this matrix and ensures that every stakeholder receives their information on schedule. When a stakeholder contacts the team asking ``what's the status?''; that is a failure of proactive communication.

\section{The SOW as a Living Document}
While the SOW is baselined at Sprint~1, it is reviewed and potentially amended at each design review:

\textbf{At PDR}: confirm that the scope, deliverables, and schedule are still realistic given the design direction selected. If the selected concept requires capabilities not in the original scope (e.g., wireless communication was not planned but the selected concept requires it), amend the SOW with client and PA approval.

\textbf{At CDR}: confirm that the build plan, budget, and timeline are still achievable. If procurement delays or design changes have shifted the schedule, amend the SOW.

\textbf{At FDR}: the SOW is the benchmark against which your Project Goal Attainment is evaluated. Any features delivered beyond the SOW scope are ``bonus''; any features not delivered require explanation (deferred to future work, design change, requirement waived).

Each SOW amendment is tracked with a revision number, date, description of change, and approver signatures.



\section{Risk Mitigation Strategy Types}
For each risk in the register, select one or more mitigation strategies:

\textbf{Avoidance}: eliminate the risk by changing the design approach. If the risk is ``motor controller IC may not be available due to supply chain shortage,'' avoid it by selecting a different, readily available IC. Avoidance is the strongest strategy but may require design compromises.

\textbf{Reduction}: decrease the likelihood or impact of the risk. If the risk is ``PCB may have a layout error,'' reduce it by: having a second team member review the layout, running DRC (Design Rule Check), and ordering 5 boards instead of 2 (so a revision still leaves spares from Rev~1). Reduction does not eliminate the risk but makes it manageable.

\textbf{Transfer}: shift the risk to a party better equipped to handle it. If the risk is ``CNC machining may exceed our skill level,'' transfer it by outsourcing to a professional machine shop (Xometry, Protolabs). The cost is higher, but the quality and schedule certainty improve.

\textbf{Acceptance}: acknowledge the risk and plan to deal with the consequences if it occurs. If the risk is ``the client may change priorities after PDR,'' accept it by: building schedule buffer, documenting the current scope in the SOW, and planning the scope change process (Chapter~7). Acceptance is appropriate for low-probability risks where the cost of mitigation exceeds the expected cost of the risk.

\textbf{Contingency planning}: for high-impact risks that cannot be fully mitigated, develop a specific contingency plan: ``If the custom motor does not arrive by Sprint~10, we will substitute [specific alternative motor] which meets [X\% of the torque requirement] and can be procured locally within 2 days.'' The contingency plan converts an open-ended crisis into a pre-planned response.



\section{What ``Done'' Looks Like: Defining Acceptance Criteria}
For every deliverable listed in the SOW, define what ``done'' means in terms the client can evaluate:

\textbf{Physical deliverables}: ``A functioning temperature control system that maintains the setpoint within $\pm$2\,\degC{} over the operating range of 20--200\,\degC{}, demonstrated in the client's presence at FDR. Delivered with: assembled prototype, firmware source code, engineering drawings, BOM, O\&M manual, and digital design archive.''

\textbf{Document deliverables}: ``A Final Design Report (FDRt) containing all sections specified in the Capstone Design Program Format Guide, formatted per program standards, submitted as PDF via Canvas and delivered to the client by email at least 3 days before the FDR presentation.''

\textbf{Presentation deliverables}: ``A 25-minute CDR presentation to the PA, client, and TAs, covering: design evolution since PDR, detailed design with engineering analysis, BOM and procurement status, manufacturing plan, VCRM status, and schedule/budget projections. Followed by 10 minutes of Q\&A.''

Vague acceptance criteria lead to disputes: the team thinks the work is done; the client expected something different. Specific criteria prevent this misalignment. When in doubt, ask: ``If I delivered exactly this, would you consider the requirement met?''

\section{The SOW Signature and Approval Process}
The SOW requires explicit approval from three parties:

\textbf{Team}: all members review and acknowledge the SOW, confirming that their roles, responsibilities, and the scope are accurately represented. Each member's agreement signifies their commitment to the project plan.

\textbf{PA}: reviews the SOW for completeness, feasibility, and alignment with capstone program requirements. The PA may request revisions before approving.

\textbf{Client}: reviews the SOW for alignment with their needs and expectations. The client's approval is documented (email confirmation or signature). After client approval, the SOW is \textbf{baselined}; changes require formal change control.

The approval process should be completed by the end of Sprint~1. Do not proceed to concept generation (Sprint~2) without an approved SOW. Building a design against an unapproved scope is building on sand.



\section{Track-Specific SOW Considerations}
While the SOW structure is the same for all tracks, the content emphasis differs:

\textbf{Build Track}: the SOW should emphasize: the physical deliverable description (what will the client receive?), the fabrication plan (where will it be built? what equipment is needed?), the procurement timeline (when must long-lead items be ordered?), and the testing strategy (where and how will the prototype be tested?). The schedule should show fabrication milestones alongside design review dates.

\textbf{Paper Track}: the SOW should emphasize: the analytical deliverable description (what models, simulations, or analyses will be produced?), the validation strategy (how will the results be verified against reference data or independent calculations?), the computational resources required (software licenses, computing time, data storage), and the documentation plan (the final deliverable is typically a detailed technical report, not a physical prototype). The schedule should show analysis milestones (model development, mesh convergence, validation, parametric studies) alongside design review dates.

\textbf{Competition Track}: the SOW should emphasize: the competition rules and compliance plan (which rules constrain the design? which are hardest to meet?), the registration and travel timeline (deadlines are typically non-negotiable), the budget allocation between build and travel, and the testing strategy (how will competition performance be validated before the event?). The schedule must account for the competition date as an absolute deadline that cannot be extended.

\textbf{Hybrid projects}: some projects combine elements of multiple tracks (e.g., a Build Track project with a significant analytical component, or a Competition Track project that also serves an external client). The SOW should clearly state which deliverables belong to each track element and how the team's effort is allocated.



\section{Common SOW Negotiation Scenarios}
These negotiations occur frequently during Sprint~1:

\textbf{``The client wants more than is achievable.''} Present a phased plan: ``In the two-semester capstone cycle, we can deliver [MVP scope]. If the project continues with a future team or through a sponsored research agreement, [Phase 2 scope] could be added.'' Most clients understand phasing once you present a concrete plan.

\textbf{``The team disagrees on the scope.''} Use the MoSCoW framework to reach consensus. If the team cannot agree on what is Must-Have vs. Should-Have, the disagreement reveals uncertainty about the client's priorities. Schedule a client meeting to clarify: ``Which of these features is essential for the system to be useful to you?''

\textbf{``The PA thinks the scope is too ambitious.''} Your PA has seen dozens of capstone projects and has a calibrated sense of what is achievable. If they say the scope is too large, take it seriously. Ask: ``Which elements would you recommend deferring?'' The PA's experience-based judgment is one of the most valuable inputs you will receive all semester.

\textbf{``The budget does not cover the desired scope.''} Present the trade-off explicitly: ``With the current budget of \$1,500, we can deliver Options~A and B. Adding Option~C requires an additional \$600. Would you prefer to request additional funding or defer Option~C?'' Never commit to scope that exceeds your budget without a documented plan for the funding gap.



\section{SOW Lessons from Past Teams}
\textbf{``We should have defined `done' more precisely.''} A team's SOW stated the deliverable as ``a working prototype.'' At FDR, the client expected a prototype with a painted enclosure, labeled connectors, and a user manual. The team delivered a functional but unfinished-looking breadboard system. Both parties were right based on their interpretation. A precise deliverable description (``a functioning prototype in a 3D-printed enclosure with labeled connectors, mounted on a demonstration base, accompanied by a printed O\&M manual'') would have prevented the misunderstanding.

\textbf{``We should have documented the client's verbal agreements.''} A client verbally approved a concept during a meeting. Three months later, at CDR, the client said ``I don't remember agreeing to that approach.'' Without meeting minutes documenting the decision, the team had no evidence. The scope discussion had to be repeated, consuming a week of Sprint~8.

\textbf{``We should have been more aggressive about Out of Scope.''} A team's SOW listed 8 In-Scope items and 0 Out-of-Scope items. By Sprint~4, the client had requested 3 additional features, each of which the team accepted because ``it wasn't explicitly excluded.'' The project became unmanageable. A team that had listed those 3 features as Out of Scope from the beginning would have had a clean basis for the scope change conversation.


\section{Practice Questions}
\begin{enumerate}[leftmargin=1.5em]
\item Write the Scope section (In Scope and Out of Scope) for a project to automate grain silo monitoring for a brewery. Include at least 3 items in each category with rationale.
\item Your client requests a new feature in Sprint~4 that was listed as ``Out of Scope'' in the SOW. Draft the email that initiates the scope change process.
\item Create three SMART goals for the first semester of a Build Track solar irrigation controller project.
\item Explain why the ``Out of Scope'' section of the SOW is as important as the ``In Scope'' section.
\end{enumerate}


